\appendix

% ---------------------------------------------------------------------
% Diapositivas de respaldo con apoyo visual reutilizando figuras.
% ---------------------------------------------------------------------
\begin{frame}{Respaldo: la maldición de la dimensionalidad}
  \begin{columns}[c]
    \column{0.5\textwidth}
      \footnotesize
      \begin{itemize}
        \item Con cosenos sobre vectores unitarios la distancia sigue siendo informativa, pero al crecer $d$ el contraste vecino cercano/lejano se desvanece~\cite{aggarwal_surprising_2001}
        \item Por eso la reducción (UMAP-10) \textbf{mejora} el agrupamiento y abarata las etapas siguientes
        \item Las métricas de calidad se calculan sobre los \textit{embeddings} \textbf{originales} (L2-normalizados) para comparar entre familias
        \item A la derecha, la relación entre la diferencia de distancia entre
            el vecino más cercano y más lejano y la distancia al más cercano,
            que tiende a cero a medida que $d$ crece.
    \end{itemize}
    \column{0.5\textwidth}
      \centering
      \includegraphics[width=\linewidth]{figures/curse_dimensionality.pdf}
  \end{columns}
\end{frame}

\begin{frame}{Respaldo: \textit{backend} de búsqueda y configuración HNSW}
  \begin{columns}[c]
    \column{0.55\textwidth}
      \centering
      \includegraphics[width=\linewidth]{figures/e6_e9_pareto.pdf}
    \column{0.45\textwidth}
      \footnotesize
      \begin{itemize}
        \item Por debajo de $N\approx1000$, pgvector exacto y HNSW son indistinguibles
        \item Por encima, HNSW se separa: su latencia se mantiene plana mientras la del exacto crece; a $N{=}6416$ es $3.2\times$ el exacto y $240\times$ SQLite
        \item E9: ninguna configuración baja de \textit{recall} $0.99$; subir \texttt{m} o \texttt{ef\_construction} llega a $1.0$ sin penalizar latencia
      \end{itemize}
  \end{columns}
\end{frame}

\begin{frame}{Respaldo: memoria pico por etapa}
  \begin{columns}[c]
    \column{0.55\textwidth}
      \centering
      \includegraphics[width=\linewidth]{figures/e7_memory.pdf}
    \column{0.45\textwidth}
      \footnotesize
      \begin{itemize}
        \item El agrupamiento \textbf{aglomerativo} materializa la matriz de distancias $O(n^2)$ y se descarta por memoria antes que por tiempo
        \item HDBSCAN, DBSCAN y OPTICS se mantienen sublineales
        \item Refuerza la elección de HDBSCAN: coste y memoria mínimos
      \end{itemize}
  \end{columns}
\end{frame}

\begin{frame}{Respaldo: ¿por qué silueta?}
  \footnotesize
  \begin{itemize}
    \item \textbf{La objeción es real}: la silueta asume clústeres compactos y
      convexos; penaliza las formas arbitrarias que producen los métodos de
      densidad (HDBSCAN). No se oculta.
    \item \textbf{Ninguna métrica se lee sola}. La comparación de \emph{extractores}
      se ancla en métricas \emph{extrínsecas} (ARI/NMI), que no asumen forma; la
      silueta es el \emph{proxy} para cuando no hay etiquetas.
    \item \textbf{Por qué la silueta como eje único}: es el único índice interno
      \emph{acotado} ($[-1,1]$) y comparable entre distintos $k$ y entre familias
      de algoritmos. Calinski--Harabasz crece con $n$ y $d$; Davies--Bouldin no
      está acotado $\Rightarrow$ malos para comparar entre ejecuciones.
    \item \textbf{A fortiori}: aun penalizada por la silueta, HDBSCAN llega a las
      mejores configuraciones $\Rightarrow$ el ranking es robusto al sesgo.
    \item El ruido ($-1$) se \textbf{excluye} de la silueta: corta en ambos
      sentidos (puede \emph{inflar} la de densidad); se reporta \texttt{noise\_ratio}
      aparte para que sea visible.
  \end{itemize}
\end{frame}

\begin{frame}{Respaldo: ¿por qué silueta? (¿sesga SupCon la métrica?)}
  \footnotesize
  \begin{itemize}
    \item SupCon optimiza compacidad intra-clase + separación: justo lo que premia
      la silueta. ¿No infla la métrica a su favor?
    \item \textbf{Distinción clave}: el ARI mide \emph{directamente} el objetivo
      (acuerdo con las etiquetas reales); la silueta es un \emph{proxy} geométrico
      de validez condicional.
    \item Optimizar hacia el objetivo es legítimo (es aprender). Sólo sería
      \enquote{Goodhart} si el proxy subiera \emph{sin} subir el objetivo (ARI).
    \item \textbf{Evidencia (E1)}: las CNN \emph{sin ajustar} lideran la silueta
      pero \emph{no} la calidad supervisada $\Rightarrow$ la silueta no favorece a
      los modelos ajustados; si acaso, lo contrario.
    \item Por eso el ranking de extractores descansa en el ARI, no en márgenes de
      silueta: no comparables entre geometrías moldeadas por objetivos distintos.
    \item (DBCV es un índice interno apto para densidad, pero sólo puntúa
      \emph{clusterers} de densidad $\Rightarrow$ no sirve de eje común con KMeans/GMM.)
  \end{itemize}
\end{frame}

\begin{frame}{Respaldo: silueta por extractor a tres escalas}
  \centering
  \includegraphics[width=0.85\linewidth]{figures/e1_silhouette_by_embedding.pdf}\\
  {\footnotesize Las CNN de ImageNet lideran la silueta \emph{intrínseca}, pero
  no la calidad \emph{supervisada} (E8a). La métrica intrínseca mide apariencia,
  no semántica de grafiti.}
\end{frame}

\begin{frame}{Respaldo: ejemplos de clústeres de estilo (E8a)}
  \begin{columns}[c]
    \column{0.33\textwidth}
      \centering
      \includegraphics[width=\linewidth]{figures/cluster_simple.pdf}\\
      {\scriptsize sencillo}
    \column{0.33\textwidth}
      \centering
      \includegraphics[width=\linewidth]{figures/cluster_elaborate.pdf}\\
      {\scriptsize elaborado}
    \column{0.33\textwidth}
      \centering
      \includegraphics[width=\linewidth]{figures/cluster_characters.pdf}\\
      {\scriptsize \textit{characters}}
  \end{columns}
  \vspace{0.4em}
  \centering {\footnotesize HDBSCAN colapsa los 4 estilos en 3 grupos coherentes:
  une \textit{tag}/\textit{throw-up} en \enquote{sencillo}.}
\end{frame}

\begin{frame}{Respaldo: autoría (E8b) bajo recuperación, no partición}
  \begin{columns}[c]
    \column{0.5\textwidth}
      \centering \scriptsize
      \setlength{\tabcolsep}{4pt}
      \begin{tabular}{lccc}
        \toprule
        \textbf{Extractor} & \textbf{MAP@5} & \textbf{Rec@5} & \textbf{MRR} \\
        \midrule
        CLIP ViT-B/32        & 0.252 & 0.268 & 0.262 \\
        DINOv2 (base)        & 0.248 & 0.268 & 0.251 \\
        ResNet50             & 0.246 & 0.236 & 0.256 \\
        \midrule
        DINOv2 autor         & 0.228 & 0.254 & 0.232 \\
        MobileNet autor      & 0.217 & 0.233 & 0.222 \\
        \bottomrule
      \end{tabular}\\[0.3em]
      {\scriptsize E8b: mejor celda por extractor (185 recortes, 87 autores).}
    \column{0.5\textwidth}
      \footnotesize
      \begin{itemize}
        \item Bajo recuperación sí hay señal: MRR $\approx 0.25$, $8\times$ el
            azar ($\approx 0.03$)
        \item Pero esa señal es \textbf{similitud genérica} (misma sesión,
            muro, color): las cabezas de autor \textbf{no superan} a los
            extractores sin ajustar, incluso quedan algo por debajo
        \item El ajuste fino entrenado \emph{solo en Cuenca} no transfiere a
            Salamanca (cambio de dominio + autores no vistos). Asimetría con
            estilo, que se entrenó mixto y se evalúa en dominio
      \end{itemize}
  \end{columns}
\end{frame}
